\documentclass{article}
\usepackage{amsmath,amssymb,amsthm}
\usepackage{algorithm}
\usepackage{algorithmic}
\usepackage{hyperref}
\usepackage{enumitem}
\setlist[description]{leftmargin=0pt}

\newtheorem{theorem}{Theorem}
\newtheorem{lemma}{Lemma}
\newtheorem{assumption}{Assumption}

\begin{document}

\section*{Gradient Descent Ascent (GDA) for Convex--Concave Saddle‑Point Problems}

\section*{Mathematical Formulation}
Let $L:\mathbb{R}^{n}\times\mathbb{R}^{m}\rightarrow\mathbb{R}$ be a \emph{convex–concave} function, i.e.
\[
L(\cdot ,y)\;\text{is convex for every }y,\qquad 
L(x,\cdot )\;\text{is concave for every }x .
\]
Define the primal variable $x\in\mathbb{R}^{n}$ and the dual variable $y\in\mathbb{R}^{m}$.  
The goal is to solve the saddle‑point problem
\[
\min_{y\in\mathbb{R}^{m}}\max_{x\in\mathbb{R}^{n}} L(x,y).
\]

Given a stepsize $\eta>0$, the Gradient Descent Ascent (GDA) iteration is
\begin{align}
x_{t+1}&=x_{t}+ \eta\,\nabla_{x} L(x_{t},y_{t}),\label{eq:gda-x}\\
y_{t+1}&=y_{t}- \eta\,\nabla_{y} L(x_{t},y_{t}),\label{eq:gda-y}
\end{align}
where $\nabla_{x}L$ and $\nabla_{y}L$ denote the partial gradients w.r.t. $x$ and $y$, respectively.
The same stepsize $\eta$ is used for both ascent (in $x$) and descent (in $y$).

\section*{Pseudocode}
\begin{algorithm}[H]
\caption{Gradient Descent Ascent (GDA)}\label{alg:gda}
\begin{algorithmic}[1]
\REQUIRE Initial point $(x_{0},y_{0})$, stepsize $\eta>0$, maximum iterations $T$.
\FOR{$t=0,\dots,T-1$}
    \STATE Compute partial gradients 
    \[
    g^{x}_{t}\gets \nabla_{x}L(x_{t},y_{t}),\qquad 
    g^{y}_{t}\gets \nabla_{y}L(x_{t},y_{t}).
    \]
    \STATE Update the primal variable (ascent)
    \[
    x_{t+1}\gets x_{t}+ \eta\, g^{x}_{t}.
    \]
    \STATE Update the dual variable (descent)
    \[
    y_{t+1}\gets y_{t}- \eta\, g^{y}_{t}.
    \]
\ENDFOR
\RETURN $(x_{T},y_{T})$ (or any averaged iterate).
\end{algorithmic}
\end{algorithm}

\section*{Dry Run Example}
Consider the scalar function 
\[
f(w)=(w-3)^{2},\qquad w\in\mathbb{R}.
\]
Although $f$ is not a saddle‑point problem, a single‑variable version of \eqref{eq:gda-x} reduces to ordinary gradient descent:
\[
w_{t+1}=w_{t}-\eta\,\nabla f(w_{t}),\qquad 
\nabla f(w)=2(w-3).
\]

\begin{description}
\item[Parameters] $w_{0}=0$, stepsize $\eta=0.1$, number of iterations $T=3$.
\end{description}

\begin{center}
\begin{tabular}{c|c|c|c}
$t$ & $w_{t}$ & $\nabla f(w_{t})$ & $f(w_{t})$\\ \hline
0 & $0$ & $2(0-3)=-6$ & $(0-3)^{2}=9$\\
1 & $0-0.1(-6)=0.6$ & $2(0.6-3)=-4.8$ & $(0.6-3)^{2}=5.76$\\
2 & $0.6-0.1(-4.8)=1.08$ & $2(1.08-3)=-3.84$ & $(1.08-3)^{2}=3.6864$\\
3 & $1.08-0.1(-3.84)=1.464$ & $2(1.464-3)=-3.072$ & $(1.464-3)^{2}=2.3665$
\end{tabular}
\end{center}
We see the objective value decreasing from $9$ to about $2.37$ after three steps.

\newpage

\section*{Assumptions}
\begin{assumption}[Convex–concave structure]\label{ass:cc}
For all $x,x'\in\mathbb{R}^{n}$ and $y,y'\in\mathbb{R}^{m}$,
\[
L(x',y)\ge L(x,y)+\langle \nabla_{x}L(x,y),x'-x\rangle,
\]
\[
L(x,y')\le L(x,y)+\langle \nabla_{y}L(x,y),y'-y\rangle .
\]
Thus $L(\cdot ,y)$ is convex and $L(x,\cdot )$ is concave.
\end{assumption}

\begin{assumption}[Smoothness]\label{ass:smooth}
$L$ has $L$‑Lipschitz continuous partial gradients, i.e. for all $z=(x,y)$ and $z'=(x',y')$,
\[
\|\nabla_{x}L(x,y)-\nabla_{x}L(x',y')\|\le L\|z-z'\|,\qquad
\|\nabla_{y}L(x,y)-\nabla_{y}L(x',y')\|\le L\|z-z'\|.
\]
Equivalently,
\[
\|\nabla L(z)-\nabla L(z')\|\le L\|z-z'\|,
\]
where $\nabla L(z)=(\nabla_{x}L,\,-\nabla_{y}L)$.
\end{assumption}

\begin{assumption}[Existence of a saddle point]\label{ass:saddle}
There exists $(x^{\star},y^{\star})$ such that
\[
L(x^{\star},y)\le L(x^{\star},y^{\star})\le L(x,y^{\star}),\qquad\forall x,y.
\]
\end{assumption}

\begin{assumption}[Stepsize]\label{ass:stepsize}
The stepsize satisfies
\[
0<\eta\le\frac{1}{L}.
\]
\end{assumption}

\section*{Main Convergence Theorem}
\begin{theorem}[Primal–dual gap of GDA]\label{thm:gap}
Let Assumptions \ref{ass:cc}–\ref{ass:stepsize} hold and denote $z_{t}=(x_{t},y_{t})$, $z^{\star}=(x^{\star},y^{\star})$.
Define the averaged iterates
\[
\bar{x}_{T}:=\frac{1}{T}\sum_{t=0}^{T-1}x_{t},\qquad
\bar{y}_{T}:=\frac{1}{T}\sum_{t=0}^{T-1}y_{t}.
\]
Then for any $T\ge1$,
\[
\underbrace{\max_{x}\,L(x,\bar{y}_{T})-\min_{y}\,L(\bar{x}_{T},y)}_{\displaystyle\text{primal–dual gap}}
\;\le\;
\frac{\|z_{0}-z^{\star}\|^{2}}{2\eta T}.
\]
Consequently, the gap converges to $0$ at the rate $O(1/T)$.
\end{theorem}

\section*{Theoretical Properties}
\begin{itemize}
\item \textbf{Stability.} The condition $\eta\le 1/L$ guarantees that the potential
$\|z_{t}-z^{\star}\|^{2}$ is non‑increasing in expectation; no explosion occurs.
\item \textbf{Hyperparameter sensitivity.}
If $\eta$ is chosen smaller than $1/L$, the bound in Theorem \ref{thm:gap} becomes looser (the denominator $\eta$ shrinks). If $\eta>1/L$, the smoothness inequality used in the proof fails and divergence may happen.
\item \textbf{Comparison to baselines.}
Standard Gradient Descent on the primal alone yields $O(1/\sqrt{T})$ for non‑strongly convex functions, while the saddle‑point structure exploited by GDA improves the rate to $O(1/T)$ under convex–concave smoothness.
\end{itemize}

\section*{Preliminaries (Definitions and Lemmas)}
\begin{lemma}[Co‑coercivity of a $L$‑smooth gradient]\label{lem:coerc}
If $\nabla L$ is $L$‑Lipschitz, then for any $z,z'\in\mathbb{R}^{n+m}$,
\[
\langle \nabla L(z)-\nabla L(z'),z-z'\rangle\ge\frac{1}{L}\|\nabla L(z)-\nabla L(z')\|^{2}.
\]
\end{lemma}
\begin{proof}
The standard proof follows from the Baillon–Haddad theorem; we omit it because it is classical.
\end{proof}

\section*{Main Proof (Step‑by‑Step Derivation)}
We prove Theorem \ref{thm:gap}.  All non‑trivial steps are numbered.

\bigskip
\noindent\textbf{Notation.}
Let $z_{t}:=(x_{t},y_{t})$, $g_{t}:=\nabla L(z_{t})=(\nabla_{x}L(x_{t},y_{t}),-\nabla_{y}L(x_{t},y_{t}))$.
The GDA update \eqref{eq:gda-x}–\eqref{eq:gda-y} can be compactly written as
\begin{equation}
z_{t+1}=z_{t}-\eta g_{t}. \label{eq:update-compact}
\end{equation}

\bigskip
\noindent\textbf{Step 1.} Expand the squared distance to the saddle point after one iteration.
\begin{align}
\|z_{t+1}-z^{\star}\|^{2}
&=\|z_{t}-\eta g_{t}-z^{\star}\|^{2} \nonumber\\
&=\|z_{t}-z^{\star}\|^{2}-2\eta\langle g_{t},z_{t}-z^{\star}\rangle+\eta^{2}\|g_{t}\|^{2}. \label{eq:expand}
\end{align}
[Step 1]

\bigskip
\noindent\textbf{Step 2.} Use convex–concave property (Assumption \ref{ass:cc}) to bound the inner product.
For any saddle point $z^{\star}=(x^{\star},y^{\star})$,
\begin{align}
\langle g_{t},z_{t}-z^{\star}\rangle
&=\langle \nabla_{x}L(x_{t},y_{t}),\,x_{t}-x^{\star}\rangle
   -\langle \nabla_{y}L(x_{t},y_{t}),\,y_{t}-y^{\star}\rangle \nonumber\\
&\ge L(x_{t},y^{\star})-L(x^{\star},y_{t}). \label{eq:innerbound}
\end{align}
[Step 2] The inequality follows from the definition of convexity in $x$ and concavity in $y$.

\bigskip
\noindent\textbf{Step 3.} Apply smoothness (Assumption \ref{ass:smooth}) to bound $\|g_{t}\|^{2}$.
Since $\nabla L$ is $L$‑Lipschitz,
\[
\|g_{t}\|^{2}= \|\nabla L(z_{t})-\nabla L(z^{\star})\|^{2}
\le L^{2}\|z_{t}-z^{\star}\|^{2}. \label{eq:smoothbound}
\]
[Step 3] (We used that $\nabla L(z^{\star})=0$ because $z^{\star}$ is a saddle point.)

\bigskip
\noindent\textbf{Step 4.} Insert \eqref{eq:innerbound} and \eqref{eq:smoothbound} into \eqref{eq:expand}.
\begin{align}
\|z_{t+1}-z^{\star}\|^{2}
&\le \|z_{t}-z^{\star}\|^{2}
      -2\eta\bigl[L(x_{t},y^{\star})-L(x^{\star},y_{t})\bigr]
      +\eta^{2}L^{2}\|z_{t}-z^{\star}\|^{2} \nonumber\\
&= \bigl(1+\eta^{2}L^{2}\bigr)\|z_{t}-z^{\star}\|^{2}
   -2\eta\bigl[L(x_{t},y^{\star})-L(x^{\star},y_{t})\bigr]. \label{eq:pre-sum}
\end{align}
[Step 4]

\bigskip
\noindent\textbf{Step 5.} Choose $\eta\le 1/L$. Then $1+\eta^{2}L^{2}\le 1+ \eta L\le 2$ and, more importantly,
\[
1+\eta^{2}L^{2}\le \frac{1}{1-\eta L}\qquad\text{(since }0\le\eta L\le1\text{)}.
\]
Multiplying both sides of \eqref{eq:pre-sum} by $(1-\eta L)$ gives
\begin{align}
(1-\eta L)\|z_{t+1}-z^{\star}\|^{2}
&\le \|z_{t}-z^{\star}\|^{2}
   -2\eta\bigl[L(x_{t},y^{\star})-L(x^{\star},y_{t})\bigr]. \label{eq:telescoping}
\end{align}
[Step 5] Rearranging terms yields a telescoping inequality.

\bigskip
\noindent\textbf{Step 6.} Sum \eqref{eq:telescoping} over $t=0,\dots,T-1$.
\begin{align}
\sum_{t=0}^{T-1}\bigl[L(x_{t},y^{\star})-L(x^{\star},y_{t})\bigr]
&\le \frac{1}{2\eta}\Bigl(\|z_{0}-z^{\star}\|^{2}
      -(1-\eta L)\|z_{T}-z^{\star}\|^{2}\Bigr) \nonumber\\
&\le \frac{\|z_{0}-z^{\star}\|^{2}}{2\eta}. \label{eq:sum-bound}
\end{align}
[Step 6] The last inequality discards the non‑negative term $(1-\eta L)\|z_{T}-z^{\star}\|^{2}$.

\bigskip
\noindent\textbf{Step 7.} Relate the summed gap to the gap of the averaged iterates.
Convexity in $x$ and concavity in $y$ imply
\[
L\bigl(\bar{x}_{T},y^{\star}\bigr)\le\frac{1}{T}\sum_{t=0}^{T-1}L(x_{t},y^{\star}),\qquad
L\bigl(x^{\star},\bar{y}_{T}\bigr)\ge\frac{1}{T}\sum_{t=0}^{T-1}L(x^{\star},y_{t}).
\]
Subtracting the second inequality from the first and using \eqref{eq:sum-bound} gives
\[
L\bigl(\bar{x}_{T},y^{\star}\bigr)-L\bigl(x^{\star},\bar{y}_{T}\bigr)
\le \frac{1}{T}\sum_{t=0}^{T-1}\bigl[L(x_{t},y^{\star})-L(x^{\star},y_{t})\bigr]
\le \frac{\|z_{0}-z^{\star}\|^{2}}{2\eta T}. \label{eq:avg-gap}
\]
[Step 7] By definition of the primal–dual gap,
\[
\max_{x}L(x,\bar{y}_{T})-\min_{y}L(\bar{x}_{T},y)
\le L(\bar{x}_{T},y^{\star})-L(x^{\star},\bar{y}_{T}),
\]
so the bound \eqref{eq:avg-gap} is exactly the claim of Theorem \ref{thm:gap}.

\bigskip
\noindent\textbf{Conclusion.} The primal–dual gap decays as $O(1/T)$, completing the proof. $\square$

\section*{Rate Derivation (With Constants)}
From \eqref{eq:avg-gap} we obtain the explicit constant
\[
C:=\frac{\|z_{0}-z^{\star}\|^{2}}{2\eta},
\qquad\text{so that}\qquad
\text{gap}_{T}\le \frac{C}{T}.
\]
If the initial distance is bounded by $D$, i.e. $\|z_{0}-z^{\star}\|\le D$, then $C\le D^{2}/(2\eta)$.

\section*{Special Cases}
\begin{description}
\item[Strong convexity / strong concavity.] If $L$ is $\mu$‑strongly convex in $x$ and $\mu$‑strongly concave in $y$, the same analysis yields a linear rate $\bigl(1-\eta\mu\bigr)^{T}$ for the distance to the saddle point.
\item[Non‑smooth case.] When only Lipschitz continuity of $L$ (but not gradient Lipschitz) holds, one can replace the plain GDA by a mirror‑proximal variant and obtain $O(1/\sqrt{T})$ rates.
\item[Stochastic gradients.] If $g_{t}$ is an unbiased estimator of $\nabla L(z_{t})$ with bounded variance, the bound becomes
\[
\mathbb{E}[\text{gap}_{T}]\le \frac{\|z_{0}-z^{\star}\|^{2}}{2\eta T}+ \frac{\eta\sigma^{2}}{2},
\]
optimizing $\eta\propto 1/\sqrt{T}$ gives $O(1/\sqrt{T})$.
\end{description}

\end{document}